\begin{frame}{Implementation of Shor's algorithm}
			\begin{figure}[H]
				\centering
\scalebox{0.70}{
\begin{quantikz}[row sep={0.7cm,between origins}, column sep=0.8cm]
        % First register (t qubits)
        \lstick{$\ket{0}_0$} & \gate{H} & \ctrl{4} & \qw & \qw & \qw & \gate[wires=4]{QFT^\dagger} & \meter{} \\
        \lstick{$\ket{0}_1$} & \gate{H} & \qw & \ctrl{3} & \qw & \qw &  & \meter{} \\
        \lstick{$\vdots$}\\
        \lstick{$\ket{0}_{t-1}$} & \gate{H} & \qw  & \qw & \qw & \ctrl{1} & \qw  & \meter{} \\
        % Second register (work register)
        \lstick{$\ket{1}$} & \qw & \gate{U^{2^{t-1}}} & \gate{U^{2^{t-2}}} & \dots   & \gate{U^{2^{0}}} & \qw & \rstick{$\ket{x^{j}º\bmod{n}}$}
\end{quantikz}
}
				\caption{Quantum circuit implementing Shor's algorithm.}
				\label{fig:shor_qc}
				
			\end{figure}


\begin{block}{Quantum Fourier Transform}
Let $\{\ket{k}\}_{k=0}^{N-1}$ be an orthonormal basis. We define the \textit{Quantum Fourier Transform} (QFT) as the unitary operation $U_{QFT}$ such that its action on the basis is,
\[U_{\text{QFT}}\ket{r}=\frac{1}{\sqrt{N}}\sum_{k=0}^{N-1}e^{i\frac{2\pi}{N}kr}\ket{k}.\]
\end{block}

\end{frame}